\begin{exbox}[固定探测面的能量相位与传播误差]
时间演化使用 $E(p)t$，固定平面读出则更方便使用 $p(E)L$。
沿向前正能支路，前章的能量归一化满足 $\int|b(E)|^2\dd E=1$。
定义该支路的理想到达振幅
\begin{equation}
 f_L(t)=\int_{m_{\rm e}c^2}^\infty\frac{\dd E}{\sqrt{2\pi\hbar}}\,
 b(E)\exp[\ii p(E)L/\hbar-\ii Et/\hbar],\qquad
 p(E)=\frac{\sqrt{E^2-m_{\rm e}^2c^4}}c.
 \label{eq:feqo-fixed-plane-energy-amplitude}
\end{equation}
Parseval 给 $\int_{\mathbb R}|f_L(t)|^2\dd t=1$，漂移只给能量振幅乘单位模相位。
这个向前支路的理想时间测量是一个 POVM；它不是对任意反射波都成立的“位置密度等于到达概率”。
真实探测器还需使用后文的效应和响应核。

为求群延迟，先直接微分而不在 $p$、波数和能量之间混用单位：
\begin{equation}
 p'(E)=\frac{E}{c\sqrt{E^2-m_{\rm e}^2c^4}}=\frac1v,\qquad
 p''(E)=-\frac{m_{\rm e}^2c^3}{(E^2-m_{\rm e}^2c^4)^{3/2}}
 =-\frac1{\gamma^3m_{\rm e}v^3}.
 \label{eq:feqo-fixed-plane-momentum-derivatives}
\end{equation}
这里 $p=\gamma m_{\rm e}v$，故第二个等式立即由
$(E^2-m_{\rm e}^2c^4)^{3/2}=c^3p^3$ 得到。
记 $\varepsilon=E-E_0$、$\tau=t-L/v_0$，除去共同载波后的振幅为
\begin{equation}
 \widetilde f_L(\tau)=\int\frac{\dd\varepsilon}{\sqrt{2\pi\hbar}}\,
 b(E_0+\varepsilon)\exp\!\left[-\frac{\ii\varepsilon\tau}{\hbar}
 -\frac{\ii L\varepsilon^2}{2\hbar\gamma_0^3m_{\rm e}v_0^3}
 +\frac{\ii L R_3(\varepsilon)}\hbar\right].
 \label{eq:feqo-fixed-plane-fresnel-kernel}
\end{equation}
其中 $|R_3|\le M_{p,3}|\varepsilon|^3/6$。
在实际支持 $|\varepsilon|\le\Delta E$ 上，丢弃三阶相位的 $L^2(\dd t)$ 振幅误差至多为
$LM_{p,3}(\Delta E)^3/(6\hbar)$；无有限支集的包还需计入尾部。
这里直接采用能量变量，不需要再把传播时间粗略替换为 $L/v_0$。
若从二阶动量 Fourier 积分的 $t=L/v_0+\tau$ 出发，另行丢弃 $\tau$ 乘二次色散的项，
还必须检查 $|\tau|\Delta p^2/(2m_\parallel\hbar)\ll1$；仅有 $v_0|\tau|\ll L$ 并不能保证该累计相位很小。

相位的驻点给局部群延迟
\begin{equation}
 t(E)=Lp'(E)=\frac L{v(E)},\qquad
 \delta t=-\frac{L\delta E}{\gamma_0^3m_{\rm e}v_0^3}+O(\delta E^2).
 \label{eq:feqo-arrival-time-shift}
\end{equation}
电子获得能量时到达更早。相邻能量梳分量的中心延迟差领先阶为
$-L\hbar\omega/(\gamma_0^3m_{\rm e}v_0^3)$。
相干源包之间是否还能干涉，需将这个分离量与源包宽度比较；
逐发中心能量随机平均的相位衰减则是另一统计操作，不能把两者统称为同一内禀退相干长度。
\end{exbox}

\begin{figure}[htbp]
 \centering
 \includegraphics[width=0.55\linewidth]{pinem_energy_ladder.pdf}
 \caption{正负整数边带的能量标记 $E_0+\ell\hbar\omega$。这张图表示能量阶数，
 不表示边带占据、真实动量等间隔或任意能区都可延伸到负无限的物理谱。}
 \label{fig:pinem-energy-ladder}
\end{figure}
